{\scriptsize
\renewcommand{\arraystretch}{1}
\setlength{\tabcolsep}{2pt}

\noindent
% ============================================================
% LEFT MINIPAGE
% ============================================================
\begin{minipage}[t]{0.495\textwidth}
\vspace{0pt}

\textbf{\large Essential Proofs II}\par
\vspace{2pt}

\begin{tabular}{
    @{}
    >{\bfseries\raggedright\arraybackslash}p{2.9cm}
    >{\raggedright\arraybackslash}p{10.5cm}
    @{}
}
\toprule

% ...

Bolzano--Weierstrass on $[a,b]$ &
A sequence $(x_n)\subseteq[a,b]$ is bounded. By Bolzano--Weierstrass,
it has a convergent subsequence
$
x_{n_k}\to x_0.
$
Since
$
a\le x_{n_k}\le b
$
for all $k$, order preservation under limits yields
$
a\le x_0\le b.
$
Thus $x_0\in[a,b]$.
\\

\midrule

Continuous on compact $\Rightarrow$ uniformly continuous &
Assume not. Then some $\varepsilon_0>0$ and sequences
$x_n,y_n\in[a,b]$ satisfy
$
|x_n-y_n|<1/n
$
but
$
|f(x_n)-f(y_n)|\ge\varepsilon_0.
$
Choose a subsequence $x_{n_k}\to x^*\in[a,b]$. Since
$
|x_{n_k}-y_{n_k}|\to0,
$
also $y_{n_k}\to x^*$. Continuity gives
$
f(x_{n_k}),f(y_{n_k})\to f(x^*),
$
hence
$
|f(x_{n_k})-f(y_{n_k})|\to0,
$
a contradiction.
\\

\midrule

Intermediate value theorem &
Assume $f(a)\le c\le f(b)$ and set
$
X:=\{x\in[a,b]\mid f(x)\le c\},
\qquad
x^*:=\sup X.
$
Choose $x_n\in X$ with $x_n\to x^*$; continuity gives
$
f(x^*)\le c.
$
If $f(x^*)<c$, continuity yields $\delta>0$ such that
$
f(y)<c
$
for $|y-x^*|<\delta$, producing $y>x^*$ in $X$, contradicting the
definition of $x^*$. Hence
$
f(x^*)=c.
$
\\

\midrule

Cauchy criterion for series &
Let
$
s_n:=\sum_{k=0}^{n}a_k.
$
Then
$
\sum_{k=0}^{\infty}a_k
$
converges iff $(s_n)$ converges, iff $(s_n)$ is Cauchy. Since
$
s_n-s_m=\sum_{k=m+1}^{n}a_k,
$
this is equivalent to
$
\forall\varepsilon>0\ \exists N\
\forall n\ge m\ge N:
\left|\sum_{k=m+1}^{n}a_k\right|<\varepsilon.
$
\\

\midrule

Absolute convergence $\Rightarrow$ convergence &
Assume
$
\sum |a_n|
$
converges. By the Cauchy criterion, for large $n\ge m$,
$
\sum_{k=m+1}^{n}|a_k|<\varepsilon.
$
Therefore
$
\left|\sum_{k=m+1}^{n}a_k\right|
\le
\sum_{k=m+1}^{n}|a_k|
<\varepsilon.
$
Thus $\sum a_n$ is Cauchy and converges. Moreover,
$
\left|\sum_{k=0}^{n}a_k\right|
\le
\sum_{k=0}^{n}|a_k|;
$
passing to the limit gives
$
\left|\sum a_k\right|\le\sum|a_k|.
$
\\

\midrule

Leibniz criterion &
Let
$
s_n:=\sum_{k=0}^{n}(-1)^ka_k
$
with $a_n\downarrow0$. Then
$
s_{2n+2}-s_{2n}=-a_{2n+1}+a_{2n+2}\le0,
$
so $(s_{2n})$ decreases, while
$
s_{2n+3}-s_{2n+1}=a_{2n+2}-a_{2n+3}\ge0,
$
so $(s_{2n+1})$ increases. Also
$
s_{2n+1}\le s_{2n}
$
and
$
s_{2n}-s_{2n+1}=a_{2n+1}\to0.
$
Thus both subsequences converge to the same limit, hence $(s_n)$ converges.
\\

\midrule

Rearrangement of an absolutely convergent series &
Let $\varphi:\mathbb N_0\to\mathbb N_0$ be bijective and
$
\sum|a_n|<\infty.
$
Given $\varepsilon>0$, choose $N$ such that
$
\sum_{n>N}|a_n|<\varepsilon.
$
Choose $M$ so that
$
\{0,\ldots,N\}
\subseteq
\{\varphi(0),\ldots,\varphi(M)\}.
$
For $m\ge M$, the difference between
$
\sum_{k=0}^{m}a_{\varphi(k)}
$
and
$
\sum_{n=0}^{\infty}a_n
$
contains only terms with index $>N$, so its absolute value is
$<\varepsilon$. Hence the rearranged series has the same sum; applying
the same argument to $|a_n|$ gives absolute convergence.
\\

\midrule

\\
\textbf{\large Solutions I}\par
\\
\midrule

Affine supremum &
Let $S\neq\varnothing$ be bounded above, $a>0$, and
$M:=\sup S$. For every $x\in S$,
$
x\le M\Rightarrow ax+b\le aM+b,
$
so $aM+b$ is an upper bound of
$
aS+b:=\{ax+b:x\in S\}.
$
Conversely, let
$
P:=\sup(aS+b).
$
Then for every $x\in S$,
$
ax+b\le P\Rightarrow x\le(P-b)/a.
$
Thus $(P-b)/a$ is an upper bound of $S$, hence
$
M\le(P-b)/a
\Rightarrow aM+b\le P.
$
Together with $P\le aM+b$,
$
\boxed{\sup_{x\in S}(ax+b)=a\sup_{x\in S}x+b}.
$
\\

\midrule

Viewing-angle maximum &
For heights $92$ and $46$ and distance $x>0$,
$
\vartheta(x)
=
\arctan\!\left(\frac{92}{x}\right)
-
\arctan\!\left(\frac{46}{x}\right).
$
Since
$
\vartheta(x)\to0
$
for $x\to0^+$ and $x\to\infty$, an interior critical point gives the
global maximum. Moreover,
$
\vartheta'(x)
=
-\frac{92}{x^2+92^2}
+
\frac{46}{x^2+46^2}.
$
Thus
$
\vartheta'(x)=0
\Longleftrightarrow
\frac{46}{x^2+46^2}
=
\frac{92}{x^2+92^2}
\Longleftrightarrow
x^2=2\cdot46^2.
$
Hence
$
\boxed{x=46\sqrt2\approx65.05}.
$
\\

\midrule

$x\sin(1/x)$ at $0$ &
Let
$
f(x)=x\sin(1/x)
$
for $x\neq0$ and $f(0)=0$. The difference quotient is
$
\frac{f(h)-f(0)}{h}
=
\sin(1/h).
$
Choose
$
h_n=\frac{2}{\pi n}\to0.
$
Then
$
\sin(1/h_n)=\sin(\pi n/2)
$
does not converge. Therefore
$
\boxed{f\text{ is not differentiable at }0}.
$
Nevertheless,
$
|f(x)|\le|x|\to0,
$
so $f$ is continuous at $0$.
\\

\midrule

$x^2\sin(1/x)$ at $0$ &
Let
$
f(x)=x^2\sin(1/x)
$
for $x\neq0$ and $f(0)=0$. Then
$
\frac{f(h)-f(0)}{h}
=
h\sin(1/h).
$
Since
$
|h\sin(1/h)|\le|h|\to0,
$
the squeeze theorem gives
$
\boxed{f'(0)=0}.
$
\\

\midrule

Piecewise exponential: corner &
Let
$
g(x)=
\begin{cases}
e^{-x},&x>0,\\
1,&x\le0.
\end{cases}
$
Then $g$ is continuous at $0$, but
$
g'_+(0)
=
\lim_{h\to0^+}\frac{e^{-h}-1}{h}
=-1,
\qquad
g'_-(0)
=
\lim_{h\to0^-}\frac{1-1}{h}
=0.
$
Since the one-sided derivatives differ,
$
\boxed{g\text{ is not differentiable at }0}.
$
\\

Cube root: vertical tangent &
For
$
f(x)=\sqrt[3]{x},
$
$
\frac{f(x)-f(0)}{x}
=
\frac{x^{1/3}}{x}
=
x^{-2/3}\to+\infty
\qquad(x\to0).
$
Hence no finite derivative exists:
$
\boxed{f\text{ is not differentiable at }0}.
$
Geometrically, the graph has a vertical tangent there.
\\

\midrule

Oscillating tangent slopes &
Let
$
h(x)=
\begin{cases}
x\sin(1/x),&x\neq0,\\
0,&x=0.
\end{cases}
$
Then
$
|h(x)-h(0)|
=
|x\sin(1/x)|
\le|x|\to0,
$
so $h$ is continuous at $0$. But
$
\frac{h(x)-h(0)}{x}
=
\sin(1/x),
$
which has no limit as $x\to0$. Thus
$
\boxed{h\text{ is continuous but not differentiable at }0}.
$
\\

\midrule

Rational--irrational function &
Let
$
f(x)=
\begin{cases}
x,&x\in\mathbb Q,\\
0,&x\notin\mathbb Q.
\end{cases}
$
At $0$,
$
|f(x)-f(0)|=|f(x)|\le|x|\to0,
$
so $f$ is continuous. Let $x_0\neq0$. Choose rational
$
r_n\to x_0
$
and irrational
$
s_n\to x_0.
$
Then
$
f(r_n)=r_n\to x_0,
\qquad
f(s_n)=0\to0.
$
The two limits differ, so
$
\boxed{f\text{ is continuous exactly at }x=0}.
$

\\

\bottomrule
\end{tabular}

\end{minipage}
\hfill
% ============================================================
% RIGHT MINIPAGE
% ============================================================
\begin{minipage}[t]{0.495\textwidth}
\vspace{0pt}

\textbf{\large Solutions II}\par
\vspace{2pt}

\begin{tabular}{
    @{}
    >{\bfseries\raggedright\arraybackslash}p{3.0cm}
    >{\raggedright\arraybackslash}p{10.4cm}
    @{}
}
\toprule

% ...

\bottomrule
\end{tabular}

\end{minipage}
}